\section{实验结果}

\subsection{主要精度对比}
表~\ref{tab:main} 在去重 DUT Anti-UAV 测试集上比较通用紧凑检测器与无人机专用检测器。SRPA-YOLO 的 Precision、Recall、AP@0.50、\mapall{} 和 AP@0.75 分别达到 95.850\%、92.201\%、95.185\%、67.071\% 和 77.329\%。由于 SRPA-YOLO 保留 YOLOv8 的单阶段训练与部署框架，YOLOv8n 是直接结构基线。相对 YOLOv8n，五项指标分别提高 3.088、8.242、5.057、6.895 和 9.750 个百分点，同时部署参数量减少 0.842 M。在统一评估协议下，Precision 与 Recall 同步提高，使“性能变化仅来自置信度工作点偏移”的解释可能性降低。

\begin{table*}[!t]
\caption{在 2199 张去重 DUT Anti-UAV 测试图像上的比较，粗体表示最优值。\label{tab:main}}
\centering\srpatablesetup
\begin{tabular}{lrrrrrr}\toprule
方法 & \shortstack{参数量\\(M)} & \shortstack{Precision\\(\%)} & \shortstack{Recall\\(\%)} & \shortstack{AP@0.50\\(\%)} & \shortstack{\mapall{}\\(\%)} & \shortstack{AP@0.75\\(\%)}\\\midrule
YOLO11n & 2.582 & 92.682 & 83.868 & 90.100 & 59.280 & 66.383\\
YOLOv8n & 3.006 & 92.762 & 83.959 & 90.128 & 60.176 & 67.579\\
EDGS-YOLOv8 & 3.091 & 94.266 & 83.244 & 90.428 & 60.347 & 66.415\\
DRBD-YOLOv8 & 3.047 & 92.658 & 83.601 & 90.203 & 60.689 & 67.459\\
YOLOv8n-P2 & 2.921 & \textbf{96.114} & 90.375 & 93.977 & 65.407 & 76.290\\
RLRD-YOLO & 2.926 & 95.781 & 90.107 & 94.467 & 65.436 & 74.757\\
\textbf{SRPA-YOLO} & \textbf{2.164} & 95.850 & \textbf{92.201} & \textbf{95.185} & \textbf{67.071} & \textbf{77.329}\\\midrule
$\Delta$ SRPA 相对 YOLOv8n & $-0.842$ & $+3.088$ & $+8.242$ & $+5.057$ & $+6.895$ & $+9.750$\\\bottomrule
\end{tabular}
\end{table*}


图~\ref{fig:comparison} 汇总各检测器的精度与紧凑性对比。

\begin{figure*}[!t]
\centering
\includegraphics[width=\textwidth]{Fig5.pdf}
\caption{去重 DUT Anti-UAV 测试集上的精度与紧凑性对比，检测指标均采用百分制。\label{fig:comparison}}
\end{figure*}

YOLOv8n-P2 用于控制高分辨率预测层带来的影响。尽管其 Precision 高 0.264 个百分点，SRPA-YOLO 的 Recall、AP@0.50、\mapall{} 和 AP@0.75 仍分别提高 1.826、1.208、1.664 和 1.039 个百分点。该对比表明，检测头任务特征分配在高分辨率尺度扩展之外仍能带来进一步增益。

图~\ref{fig:complexity} 展示部署参数量与 \mapall{} 的关系。SRPA-YOLO 位于左上区域，在对比模型中同时获得最高 \mapall{} 和最小参数量。该工作点与“使用紧凑共享字典，并仅向分类分配附加容量”的设计目标一致。

\begin{figure}[!t]
\centering
\includegraphics[width=\columnwidth]{Fig6.pdf}
\caption{模型复杂度权衡。横轴为部署参数量，纵轴为 \mapall{}，箭头表示更优方向。\label{fig:complexity}}
\end{figure}

\subsection{共享/私有秩与机制消融}
表~\ref{tab:rank} 在验证集上独立分析共享秩与私有秩。共享秩从 48 增至 64 或 80 时，\mapall{} 并未单调增加，说明统一增宽不足以确定最优配置。48+16 取得最高验证 \mapall{}（60.874\%），64+16 则取得最高 AP@0.75（69.986\%）、AP@0.50（91.850\%）和 Precision（96.615\%）。依据式~\eqref{eq:selection}，两种分配均处于最高 \mapall{} 的 0.3 个百分点内，64+16 因具有更强的高 IoU 定位性能而被选用。

\begin{table*}[!t]
\caption{共享秩与私有秩的验证集消融。\label{tab:rank}}
\centering\srpatablesetup
\begin{tabular}{rrrrrrr}\toprule
$r_s$ & $r_p$ & \shortstack{参数量\\(M)} & \shortstack{P\\(\%)} & \shortstack{R\\(\%)} & \shortstack{AP@0.50\\(\%)} & \shortstack{\mapall{}/\\AP@0.75 (\%)}\\\midrule
48 & 0  & 2.093 & 96.089 & 87.182 & 91.667 & 60.508 / 69.337\\
64 & 0  & 2.128 & 96.378 & 86.570 & 91.703 & 60.454 / 69.754\\
80 & 0  & 2.164 & 96.073 & \textbf{87.743} & 91.569 & 60.513 / 69.262\\
48 & 16 & 2.128 & 96.135 & 87.310 & 91.798 & \textbf{60.874} / 69.618\\
\textbf{64} & \textbf{16} & \textbf{2.164} & \textbf{96.615} & 86.570 & \textbf{91.850} & 60.646 / \textbf{69.986}\\\bottomrule
\end{tabular}
\end{table*}


表~\ref{tab:mechanism} 展示从完整非线性检测头到最终 SRPA 配置的累积变化。加入密集 DFL 蒸馏与检测头压缩精度的恢复相关，随后增加的置信度聚焦项对应 0.384 个 \mapall{} 百分点增量。在累积协议下，迁移初始化对应 Recall 和 \mapall{} 分别提高 2.660 和 0.678 个百分点；共享表征增至 64 通道后，\mapall{} 和 AP@0.75 分别增加 0.531 和 1.674 个百分点。最终，16 通道分类私有残差使 \mapall{} 达到 67.071\%，AP@0.75 达到 77.329\%。各行均保留前序组件，并报告相应累积配置的性能。共享 80 通道对照获得相近的 67.121\% \mapall{}，而 64+16 分配还满足式~\eqref{eq:jacobian} 的零边界框雅可比，并由预先定义的验证集定位准则选出。结合秩消融，结果支持受任务约束的容量分配优于不受约束的检测头增宽。

\begin{table*}[!t]
\caption{去重 DUT Anti-UAV 测试集上的机制消融。\label{tab:mechanism}}
\centering\srpatablesetup
\begin{tabular}{p{0.27\textwidth}rrrrrr}\toprule
配置 & \shortstack{参数量\\(M)} & \shortstack{P\\(\%)} & \shortstack{R\\(\%)} & \shortstack{AP@0.50\\(\%)} & \shortstack{\mapall{}\\(\%)} & \shortstack{AP@0.75\\(\%)}\\\midrule
完整非线性头 & 2.603 & 95.922 & 90.954 & 94.736 & 66.046 & 75.503\\
Shared-48 + 密集 DFL KD & 2.093 & 96.152 & 90.954 & 94.836 & 65.363 & 74.929\\
$+$ 聚焦定位 KD & 2.093 & 97.084 & 90.285 & 95.018 & 65.747 & 75.573\\
$+$ 高容量初始化迁移 & 2.093 & 96.484 & 92.945 & \textbf{95.493} & 66.425 & 75.472\\
$+$ 共享秩 64 & 2.128 & 96.122 & 92.157 & 95.366 & 66.956 & 77.146\\
$+$ 私有秩 16 & 2.164 & 95.850 & 92.201 & 95.185 & \textbf{67.071} & \textbf{77.329}\\\bottomrule
\end{tabular}
\end{table*}


\subsection{定位与尺度特性}
图~\ref{fig:scaleiou} 在 0.50--0.95 的 IoU 阈值范围内直接比较 SRPA-YOLO 与 YOLOv8n。SRPA-YOLO 在每个阈值上均取得更高 AP：IoU 0.50 时差值为 5.057 个百分点，IoU 0.75 时增至 9.750 个百分点，IoU 0.85 时仍为 8.803 个百分点。严格重叠条件下仍保持优势，说明实测改进包含定位因素，而不局限于 IoU 0.50 下的分类置信度变化。

\begin{figure*}[!t]
\centering
\includegraphics[width=0.96\textwidth]{Fig7.pdf}
\caption{SRPA-YOLO 与 YOLOv8n 的定位对比。左：各 IoU 阈值下的 AP；右：综合检测指标。\label{fig:scaleiou}}
\end{figure*}

相对 YOLOv8n，最大的综合提升出现在 Recall（+8.242 个百分点），同时 AP@0.50 和 AP@0.75 均增加。在远距离无人机监测中，更高 Recall 可直接减少漏检，而 Precision 的同步提升限制了附加误警。P2 预测层为 tiny 无人机提供空间支撑，SRPA 检测头则使分类私有细化不进入边界框回归输入。该结构解释与低、高 IoU 阈值下的实测增益一致。

\subsection{定性检测分析}
图~\ref{fig:detectioncomparison} 在置信度阈值为 0.25 时，对相同 tiny UAV 实例上的 YOLOv8n 与 SRPA-YOLO 检测结果进行比较。当预测框与真值框的 IoU 不低于 0.50 时，将其判定为正确检测。三个所选样例中，YOLOv8n 均未产生满足匹配条件的预测，而 SRPA-YOLO 能够恢复目标。这些代表性案例从视觉层面印证了测试集定量指标所反映的性能提升；局部放大图展示了所恢复预测与仅占少量像素目标的空间对齐情况。

\begin{figure*}[!t]
\centering
\includegraphics[width=\textwidth]{Fig8.pdf}
\caption{去重 DUT Anti-UAV 测试集上的同图检测对比。黄色虚线框表示真值，红色实线框表示 SRPA-YOLO 预测；“Missed”表示 YOLOv8n 不存在同时满足置信度 $\geq0.25$ 和 IoU $\geq0.50$ 的预测。\label{fig:detectioncomparison}}
\end{figure*}

图~\ref{fig:tinydetections} 进一步展示 SRPA-YOLO 在天空、建筑、植被和低对比度背景中的成功检测结果。所选目标经 640 像素 letterbox 变换后的最长边仅为 5.7--6.3 px。尽管目标空间支撑有限且背景占比很高，保留的 checkpoint 在各所选场景中均形成了定位响应。这些案例直观展示了模型在天空、植被、建筑和低对比度背景中的小目标检测能力，并与前述测试集指标保持一致。

\begin{figure*}[!t]
\centering
\includegraphics[width=0.90\textwidth]{Fig9.pdf}
\caption{SRPA-YOLO 对 tiny UAV 的成功检测结果。每组右侧区域为对应原图中的目标局部放大，红色框和标签来自最终保留 checkpoint 的预测。\label{fig:tinydetections}}
\end{figure*}

\subsection{匹配随机种子稳定性}
表~\ref{tab:seeds} 汇总了阶段 II 随机种子 40、41 和 42 的实验结果。各实验采用相同的阶段 I 共享基底，并独立初始化分类私有参数。SRPA-YOLO 的 Recall、AP@0.50 和 \mapall{} 标准差分别为 0.112、0.022 和 0.082 个百分点，表明分类私有细化在不同初始化条件下保持了较低的性能波动。相对 RLRD-YOLO，Recall、\mapall{} 和 AP@0.75 的平均配对提升分别为 2.709、1.881 和 2.150 个百分点，说明所提出的分类私有表征能够在多个匹配种子下持续改善检测精度。

\begin{table*}[!t]
\caption{三个匹配阶段 II 随机种子在去重 DUT Anti-UAV 测试集上的均值与样本标准差（\%，$n=3$）。各实验使用相同的阶段 I 共享基底，仅独立初始化私有参数。\label{tab:seeds}}
\centering\srpatablesetup
\begin{tabular}{lrrrrr}\toprule
方法 & \shortstack{Precision\\(\%)} & \shortstack{Recall\\(\%)} & \shortstack{AP@0.50\\(\%)} & \shortstack{\mapall{}\\(\%)} & \shortstack{AP@0.75\\(\%)}\\\midrule
RLRD-YOLO & $95.964\pm0.766$ & $89.463\pm0.785$ & $94.661\pm0.106$ & $65.199\pm0.639$ & $75.196\pm0.564$\\
SRPA-YOLO & $96.167\pm0.031$ & $92.172\pm0.112$ & $95.304\pm0.022$ & $67.080\pm0.082$ & $77.345\pm0.233$\\
配对差值 & $+0.204\pm0.736$ & $+2.709\pm0.773$ & $+0.644\pm0.109$ & $+1.881\pm0.721$ & $+2.150\pm0.332$\\\bottomrule
\end{tabular}
\end{table*}


\subsection{延迟与部署复杂度}
训练图约含 2.18 M 参数，在 640 像素输入下为 8.5 GFLOPs。结构融合移除训练期状态后，部署 checkpoint 含 2.164 M 参数。表~\ref{tab:latency} 给出同进程配对测试，其中每个检测器均与同场 YOLOv8n 对照共同测量；表~\ref{tab:cleanlat} 给出独立低负载最终验收结果。

\begin{table*}[!t]\caption{同进程配对延迟。batch 1、640、FP16、融合前向+NMS、80 次预热、300 次计时；Control 为同步 YOLOv8n。\label{tab:latency}}\centering\srpatablesetup
\begin{tabular}{lrrrrrr}\toprule
&\multicolumn{3}{c}{第1轮}&\multicolumn{3}{c}{第2轮}\\\cmidrule(lr){2-4}\cmidrule(lr){5-7}
方法&ms&FPS&Control ms&ms&FPS&Control ms\\\midrule
YOLO11n&21.6357&46.22&16.8632&19.1505&52.22&15.0139\\
YOLOv8n-P2&19.1953&52.10&15.9227&16.9629&58.95&14.1962\\
EDGS*&19.5146&51.24&14.6088&22.0043&45.45&16.3604\\
DRBD*&29.0158&34.46&14.5733&31.7115&31.53&16.8584\\
RLRD-YOLO&22.5644&44.32&16.0771&22.7773&43.90&16.2647\\
\textbf{SRPA-YOLO}&\textbf{16.2469}&\textbf{61.55}&16.4922&\textbf{15.0936}&\textbf{66.25}&15.5909\\\bottomrule
\end{tabular}\end{table*}
\begin{table*}[!t]\caption{SRPA-YOLO 与同场 YOLOv8n 对照在相同 batch 1、640 像素、FP16、融合前向+NMS 协议下的独立低负载最终验收结果。\label{tab:cleanlat}}\centering\srpatablesetup
\begin{tabular}{lrrrr}\toprule
轮次&SRPA-YOLO(ms)&SRPA-YOLO(FPS)&YOLOv8n(ms)&YOLOv8n(FPS)\\\midrule
1&9.7435&102.63&9.8859&101.15\\2&9.8574&101.45&9.9893&100.11\\\bottomrule
\end{tabular}\end{table*}


在另一组低负载验收测试中，SRPA-YOLO 两轮延迟为 9.7435 和 9.8574 ms（102.63 和 101.45 FPS），同场测量的 YOLOv8n 对照为 9.8859 和 9.9893 ms（101.15 和 100.11 FPS）。二者两轮均值分别为 9.8005 和 9.9376 ms。在所报告的配对测试协议下，SRPA-YOLO 在保持相当帧时间的同时，将部署参数量减少了 28.0\%。

图~\ref{fig:speedtradeoff} 展示所报告的两轮平均延迟与主要精度指标之间的关系。在已测模型和当前硬件协议范围内，SRPA-YOLO 位于有利的精度--速度边界：精度相近的模型具有更高实测延迟，而更快的模型未达到相同 \mapall{}。

\begin{figure}[!t]
\centering
\includegraphics[width=0.78\columnwidth]{Fig10.pdf}
\caption{基于两轮平均延迟与 \mapall{} 的精度--速度权衡，SRPA-YOLO 位于有利 Pareto 区域。\label{fig:speedtradeoff}}
\end{figure}
